\scanpage{284}
\mspage{- 32 -}\hfill IX
% archive sheet number: 284
% mounting strip and printed label: 350

\uline{Proposition} 3.5. (\uline{critère galoisien de réduction semi-stable}). \uline{On suppose} $S$ \uline{un trait}. \uline{Les conditions suivantes sur le schéma abélien} $A_K$ \uline{sont équivalentes} (\uline{les notations étant celles de} (2.5.3), \uline{appliquées à la situation déduite de} $(A_K,S)$ \uline{par passage au hensélisé de} $S$) :

(i) $A_K$ \uline{a réduction semi-stable sur} $S$ (3.4).

(ii) \uline{On a}, \uline{avec les notations} (2.5.3), (2.5.4) :
\[
V\handwritten{\subset}V^{\handwritten{\perp}}\quad .
\]

(iii) \uline{On a pour la partie torique} $W$ \uline{du module de Tate} :
\[
W=V^{\handwritten{\perp}}\quad .
\]

(iv) \uline{L'action du groupe d'inertie} $I$ \uline{sur} $U=T_{\handwritten{\ell}}(A(\overline{K}))$ \uline{est unipotente d'échelon} 2 (\uline{i}.\uline{e}. \uline{il existe un sous-module} $U'$ \uline{de} $U$ \uline{stable par} $I$ \uline{tel que} $I$ \uline{opère trivialement sur} $U'$ \uline{et sur} $U/U'$).

On peut supposer $S$ hensélien (3.4.0). Avec les notations de 2.1.9, (i) signifie que l'on a
\[
\handwritten{\lambda}=0\quad ,
\]
ce qui en vertu de (2.5.5) équivaut à $V=V\handwritten{\cap}V^{\handwritten{\perp}}$ i.e. à (ii). Comme $W=V\handwritten{\cap}V^{\handwritten{\perp}}$ par le théorème d'orthogonalité 2.4, on a (ii) $\handwritten{\Longleftrightarrow}$ (iii). D'ailleurs (iv) signifie aussi que $I$ opère trivialement sur $U^{I}=V$, ce qui par dualité signifie qu'il opère trivialement sur $V^{\handwritten{\perp}}$ i.e. que $V^{\handwritten{\perp}}\handwritten{\subset} U^{I}=V$, ce qui est (ii). Cela achève la démonstration.

\uline{Remarques} 3.5.1. Les conditions (ii) et (iv) ne dépendent visiblement que de l'action de $I$ sur le $\mathbb{Q}_{\handwritten{\ell}}$-vectoriel $U\otimes_{\mathbb{Z}_{\handwritten{\ell}}}\mathbb{Q}_{\handwritten{\ell}}$, et auraient pu être formulées en termes de cette action. Signalons aussi que ces conditions sont manifestement stables par tout changement de base par un morphisme
